% Appendix B: secondary derivations. Everything here is proved in the text or is a CPU
% finding cited from the evidence register; the main text states what each result is for.
\section{Secondary derivations}\label{app:derivations}
\subsection{The directional margin}\label{app:directional}
When the information is about the hidden state rather than the logits, a box of logit bounds is loose; the certificate below uses the direction of each competitor instead.
\begin{proposition}[Directional margin]\label{prop:directional}
Let the logits be $z=Wh$ in real arithmetic and suppose $\norm{h-\tilde h}_2\le\varepsilon_h$ for a known $\tilde h$. Then $c$ is the strict argmax of $Wh$ for every such $h$ if and only if
\[(w_c-w_j)^\top\tilde h>\norm{w_c-w_j}_2\,\varepsilon_h\quad\text{for all }j\neq c.\]
This condition is implied by, and can be much weaker than, the box condition obtained from the row bounds $|z_j-w_j^\top\tilde h|\le\norm{w_j}_2\varepsilon_h$.
\end{proposition}
\begin{proof}
For fixed $j$, $\min_{\norm e_2\le\varepsilon_h}(w_c-w_j)^\top(\tilde h+e)=(w_c-w_j)^\top\tilde h-\norm{w_c-w_j}_2\varepsilon_h$ by Cauchy--Schwarz, with equality at $e=-\varepsilon_h(w_c-w_j)/\norm{w_c-w_j}_2$. So $z_c>z_j$ throughout the ball exactly when the stated inequality holds. The box condition reads $(w_c-w_j)^\top\tilde h>(\norm{w_c}_2+\norm{w_j}_2)\varepsilon_h$, and $\norm{w_c-w_j}_2\le\norm{w_c}_2+\norm{w_j}_2$.
\end{proof}
The directional certificate is not free: it needs $\norm{w_c-w_j}_2$ for every competitor, which for all $j$ costs as much as a head pass. It is useful after a box screen has reduced the competitors to a few rows, and for the stock contract the gap condition of Theorem~\ref{thm:stock} then applies on top of it, as it does for the refinement ladder. It is also a real-arithmetic statement; Section~\ref{sec:selfevidence} supplies the floating-point terms that turn such statements into enclosures of computed values. Section~\ref{sec:tail} uses it for the certified-tail kill test.


\subsection{Computed norms and outward rounding}\label{app:outward}
The envelope of Theorem~\ref{thm:envelope} is a real number; the kernel evaluates it in floating point. The two lemmas below make the evaluated interval an enclosure: the first computes norms of $h$ that are never below the exact norms, and the second rounds the endpoints outward.

\begin{lemma}[Computed norms]\label{lem:norms}
In FP32 with round to nearest and $D\le2^{14}$, let $\hat\sigma$ be any computed sum of the rounded squares $\fl(h_j^2)$ with depth $d\le2^{14}$, $\hat r\ge(1-2^{-16})\sqrt{\hat\sigma}$ any square root, and $N_2=\fl(\fl(\hat r(1+2^{-10}))+2^{-50})$. Then $N_2\ge\norm h_2$. Likewise, if $\hat\sigma_1$ is any computed sum of the $|h_j|$ with depth $d\le2^{13}$, then $N_1=\fl(\fl(\hat\sigma_1(1+2^{-10}))+2^{-50})\ge\norm h_1$.
\end{lemma}
\begin{proof}
A rounded nonnegative value, and a rounded addition of nonnegative values, returns at least $(1-u)$ times its exact value less $\mu$ (which also covers flush-to-zero). For $N_2$ the squares add one rounding to the $d$ levels of the sum, so $\norm h_2^2\le(1-u)^{-(d+1)}\hat\sigma+3D\mu$ and $\norm h_2\le(1-u)^{-(d+1)/2}(1-2^{-16})^{-1}\hat r+2^{-54}$ for $D\le2^{14}$. The computed $N_2\ge(1-u)^2(1+2^{-10})\hat r+2^{-51}$, and $(1-u)^2(1+2^{-10})\ge(1-u)^{-(d+1)/2}(1-2^{-16})^{-1}$ because $(d+1)u/2\le2^{-11}+2^{-25}$. For $N_1$ the terms $|h_j|$ are exact, so $\norm h_1\le(1-u)^{-d}(\hat\sigma_1+D\mu)$, and $(1-u)^2(1+2^{-10})\ge(1-u)^{-d}$ holds for $d\le2^{13}$ (then $du\le2^{-11}$) but fails at $d=2^{14}$, where the margin is about $-6\times10^{-7}$; there $N_1$ must be inflated by $1+2^{-9}$ instead.
\end{proof}

\begin{lemma}[Outward rounding of the computed interval]\label{lem:outward}
In a format with $p\ge12$ and round to nearest, let $\tilde z$ be representable, let $\beta^\ast$ be a real bound with $|z-\tilde z|\le\beta^\ast$, and let $b$ be computed from representable upper bounds by at most $k\le59$ rounded operations ($+$, $\times$, $\sqrt{\ }$) on nonnegative values along any path of the expression, so that $b\ge(1-u)^k\beta^\ast-k\mu$. Set
\[\hat b=\fl\bigl(\fl(\fl(b(1+2^{6-p}))+\fl(2^{4-p}|\tilde z|))+2^{e_{\min}+26}\bigr),\quad lo=\fl(\tilde z-\hat b),\quad hi=\fl(\tilde z+\hat b).\]
Then $lo\le z\le hi$, so every comparison between computed endpoints is sound.
\end{lemma}
\begin{proof}
Each rounding of a nonnegative value loses at most a factor $(1-u)$ and $\mu$, so $\hat b(1-u)\ge(1-u)^4(1+64u)b+16u(1-u)^3|\tilde z|+(1-u)^22^{e_{\min}+26}-4\mu$. Also $lo\le\tilde z-\hat b(1-u)+u|\tilde z|+\mu$. It suffices that $\hat b(1-u)\ge\beta^\ast+u|\tilde z|+\mu$: the $|\tilde z|$ terms satisfy $16u(1-u)^3\ge u$, the $b$ terms satisfy $(1-u)^{4+k}(1+64u)\ge(1-63u)(1+64u)\ge1$ for $u\le2^{-12}$, and $2^{26}\mu$ absorbs the absolute terms. The upper endpoint is symmetric.
\end{proof}
For FP32 the constants are $1+2^{-18}$, $2^{-20}|\tilde z|$ and $2^{-100}$.

\subsection{A summary that supports both selection and mass}\label{app:summary}
For a set of labelled logits, keep its top $k$ items under descending score then ascending token index, and the mass of every discarded item. Two such summaries merge by taking the top $k$ of their retained union and adding the newly discarded mass to the two old tails.

\begin{proposition}[Top-$k$ plus tail is a compositional summary]\label{prop:summary}
For disjoint finite domains the merge produces the same top-$k$ items and tail mass as summarizing the union, and it is associative and commutative in exact arithmetic under the common total order.
\end{proposition}
\begin{proof}
An item omitted from a local top-$k$ has at least $k$ locally superior items and cannot enter the union's top-$k$, so every possible winner is retained. The original tails and the newly discarded items partition the final tail, so summing them gives its mass without subtracting nearly equal totals. Every parenthesization and order represents the same union.
\end{proof}
This is a sufficient-statistic construction from familiar reductions, and floating-point tail masses are not associative. Its value is an interface: one epilogue can serve candidate generation, bounds and completion. Retaining a small set $R_c$ per tile also tightens the mass bounds: recompute its target logits exactly and transport only the rest,
\begin{equation}
\sum_{i\in R_c}e^{z_i^t}+e^{a_c-\epsilon_c}Z^d_{c\setminus R_c}\le Z_c^t\le\sum_{i\in R_c}e^{z_i^t}+e^{a_c+\epsilon_c}Z^d_{c\setminus R_c},\label{eq:tail}
\end{equation}
which is never looser than transporting the whole tile, at the cost of $|R_c|$ extra row projections.

\subsection{How much of a sample can be emitted before the scores are known}\label{sec:commonmass}
A race needs the winner's interval to separate. A different construction can emit a token with some probability before any candidate is separated, and then finish the rest of the law after more computation. Let $\mathcal P(I)$ be the set of target laws compatible with the information $I$ held so far.

\begin{proposition}[Common mass]\label{prop:commonmass}
Consider procedures that, before acquiring more information, emit token $i$ with probability $r_i$ (the same for every $p\in\mathcal P(I)$) and otherwise continue until $p$ is known. Such a procedure outputs exactly $p$ for every $p\in\mathcal P(I)$ if and only if $r_i\le\inf_{p\in\mathcal P(I)}p_i$ for every $i$, and then it may finish by sampling the residual $(p-r)/(1-\alpha)$ with $\alpha=\sum_ir_i$. The largest immediate probability is
\[\alpha^\ast(I)=\sum_i\inf_{p\in\mathcal P(I)}p_i.\]
\end{proposition}
\begin{proof}
If the procedure is exact, the final probability of $i$ is at least the probability $r_i$ of emitting it immediately, for every compatible $p$; so $r_i\le\inf_pp_i$. Conversely, if $r_i\le\inf_pp_i$ then $p-r\ge0$ and sums to $1-\alpha$ for every compatible $p$, so the residual is a law, and the output law is $r+(1-\alpha)\cdot(p-r)/(1-\alpha)=p$.
\end{proof}
For two candidate laws $\mathcal P=\{p,q\}$, $\alpha^\ast=\sum_i\min(p_i,q_i)=1-\TV(p,q)$, the overlap of a maximal coupling, which is exactly the acceptance probability of speculative sampling~\cite{specdecode,chen2023specsampling}. The proposition says the same idea applies to any set of laws the evidence leaves open, and that nothing can do better without more information.

A two-token example shows that exact sampling can finish where greedy selection cannot. Suppose the evidence only says that the probability of token~1 lies in $[0.49,0.51]$. Then $\alpha^\ast=0.49+0.49=0.98$. Draw $U$ uniform on $[0,1)$: emit token~1 if $U<0.49$, token~0 if $U\ge0.51$, and otherwise compute $p$ and emit token~1 exactly when $U<p$, reusing the same $U$. The output has law $p$, and 98\% of draws never need the scores; the greedy decision, by contrast, is undecided because the interval contains $1/2$. For two tokens this procedure is inverse-CDF sampling with bounds on the CDF. In general it gives equality in distribution, not seeded equality with a dense race.

\begin{proposition}[Bounded logit error]\label{prop:boundedrange}
Let $p=\softmax(z/T)$ and $q=\softmax(\tilde z/T)$, and suppose $\max_i(z_i-\tilde z_i)-\min_i(z_i-\tilde z_i)\le\Delta$. Then $p_i\ge e^{-\Delta/T}q_i$ for every $i$. Consequently, with $\alpha=e^{-\Delta/T}$:
\begin{enumerate}
\item sampling $q$ with probability $\alpha$, and otherwise computing $p$ and sampling $(p-\alpha q)/(1-\alpha)$, samples $p$ exactly;
\item always sampling $q$ gives $\TV(p,q)\le1-\alpha$.
\end{enumerate}
\end{proposition}
\begin{proof}
Write $d_i=z_i-\tilde z_i$ with $m\le d_i\le M$ and $M-m\le\Delta$. Then $p_i/q_i=e^{d_i/T}/\sum_jq_je^{d_j/T}\ge e^{m/T}/e^{M/T}\ge e^{-\Delta/T}$. Part~1 is Proposition~\ref{prop:commonmass} with $r=\alpha q$. For part~2, $\TV(p,q)=\sum_i(q_i-p_i)_+\le\sum_i(1-\alpha)q_i=1-\alpha$.
\end{proof}
With self-evidence the range is at most twice the largest envelope half-width, so $\alpha\ge e^{-2\max_i\beta_i/T}$. The trade-off against the race is different: this sampler needs the cheap pass's full softmax and, with probability $1-\alpha$, the whole exact head, whereas the race re-scores a candidate set whose expected size grows like $e^{4\beta/T}$ and returns the seeded token.

\subsection{Acceptance guards and residual races}\label{sec:guards}\label{app:guards}
For one reached proposal $X=x$ write $w_x=e^{z_x}$ and $Z=\sum_ie^{z_i}$, and draw $U\in[0,1)$ with the reference convention. For $q_x>0$, acceptance is equivalent to
\begin{equation}Uq_xZ<w_x,\label{eq:accept}\end{equation}
and the $\min(1,\cdot)$ need not be evaluated because $U<1$ already handles $p_x\ge q_x$.

\begin{theorem}[Acceptance guards with enclosures]\label{thm:guards}\label{thm:guard}
Let $w_x\in[N^-,N^+]$ and $Z\in[Z^-,Z^+]$ with $N^->0$, let $0<q_x\le1$ be the actual proposal mass and $U\in[0,1)$. Then $Uq_xZ^+<N^-$ implies acceptance, $Uq_xZ^-\ge N^+$ implies rejection, the two never both hold, and replacing an enclosure by a tighter one cannot reverse a certified decision. For $U$ uniform and independent of the enclosures, neither guard fires with probability
\[P_?=\min\bigl(1,N^+/(q_xZ^-)\bigr)-\min\bigl(1,N^-/(q_xZ^+)\bigr).\]
If $N^-=N^+=w_x$ and every row outside a re-scored set $R$ has a mass enclosure whose upper and lower ends differ by a factor at most $\kappa$, then with $\pi=\sum_{i\notin R}p_i$,
\[P_?\le\frac{p_x}{q_x}\cdot\frac{(\kappa-1)\pi}{1-\pi(1-1/\kappa)}.\]
\end{theorem}
\begin{proof}
The implications are transitivity with~\eqref{eq:accept}; both firing would give $N^+\le Uq_xZ^-\le Uq_xZ^+<N^-$. A tighter enclosure stays on the same side of the true values. The unresolved uniforms form $[N^-/(q_xZ^+),N^+/(q_xZ^-))\cap[0,1)$. For the bound, $P_?\le(w_x/q_x)(Z^+-Z^-)/(Z^+Z^-)$; $Z^+-Z^-\le\sum_{i\notin R}e^{z_i}(\kappa-1)$ because each lower mass is at most $e^{z_i}$; and $Z^-\ge Z-\sum_{i\notin R}e^{z_i}(1-1/\kappa)$ because each lower mass is at least $e^{z_i}/\kappa$. Divide by $Z$.
\end{proof}
These guards are analysis in this revision: the engine path adopted in Section~\ref{sec:fixednoise} uses fixed-noise verification, which needs no partition function, and a certified version of the stock rule would still pay for a dense residual row. Transported tile masses (Theorem~\ref{thm:transport}) and exponentiated self-evidence intervals (Theorem~\ref{thm:envelope}) both supply $Z^\pm$. With self-evidence $\kappa=e^{2\hat\beta}$ times the outward rounding of the exponential; at $\beta=1$ ($\kappa\approx7.4$), keeping $P_?$ below 10\% at $p_x/q_x\approx1$ requires the unre-scored mass $\pi$ to be below about 1.5\%. After an adaptive refinement chosen using $U$, the probability statement no longer applies, although the guards remain sound pointwise. The Lean formalization proves the guards with enclosed numerators and their exclusivity (\code{formal/CertifiedArgmax.lean}).

\paragraph{Residual and bonus work}\label{sec:residual}
A rejected proposal still requires a sample from $(p-q)_+$, and knowing that rejection occurred does not supply one. When all drafted tokens are accepted, a bonus token needs a target sample. A certificate that decides acceptance cheaply and then completes the whole head for the residual saves little; Appendix~\ref{app:synthetic} shows this happening. Two facts about the engine and one construction reduce the problem.

First, SGLang's default sampled verification uses point-mass proposals (Section~\ref{sec:stock}). With $q$ a point mass at $x$, the residual is the target restricted to the untried tokens, a race over $V\setminus\{x\}$ that needs no partition function. Second, under fixed-noise verification there is no residual at all. Third, for a proposal with small support there is an exact bounded race.

Let $Q=\{i:q_i>0\}$. With unnormalized target weights $w_i$ and partition $Z$, the residual can be sampled from weights
\begin{equation}
r_i=(w_i-Zq_i)_+,\qquad\frac{r_i}{\sum_jr_j}=\frac{(p_i-q_i)_+}{\sum_j(p_j-q_j)_+},\label{eq:residualweights}
\end{equation}
and $r_i=w_i$ outside $Q$. Given $Z^-\le Z\le Z^+$, positive-part monotonicity gives $(w_i-Z^+q_i)_+\le r_i\le(w_i-Z^-q_i)_+$ inside $Q$.

\begin{theorem}[Bounded residual race]\label{thm:residualrace}
Assume valid enclosures of the target weights and of $Z$, the actual proposal law, a nonzero residual, and a priority field $\pi_i=e^{G_i}>0$ that is fixed and independent of the proposal draw, the acceptance uniform and the rejection event. A refinement procedure that returns only an interval-certified ordered maximizer of $r_i\pi_i$, or completes the dense race, returns the same token as the dense race. With i.i.d.\ standard Gumbel $G_i$ it samples the residual law~\eqref{eq:residualweights} exactly.
\end{theorem}
\begin{proof}
The enclosures of $r_i$ multiplied by $\pi_i>0$ enclose the priorities. A certified maximizer dominates every possible competitor, including the actual one, and a completed evaluation yields the dense ordered maximizer. The distributional statement is Gumbel-max applied to the nonzero residual weights.
\end{proof}

\begin{theorem}[Residual races with self-evidence]\label{thm:residualself}
Replace the enclosures above by $r_i\in[(N^-_i-Z^+q_i)_+,(N^+_i-Z^-q_i)_+]$ for $q_i>0$ and $r_i\in[N_i^-,N_i^+]$ for $q_i=0$. The conclusions of Theorem~\ref{thm:residualrace} hold. For a point-mass proposal at $x$, $r_x=0$ as soon as $Z^-\ge N^+_x$, and the residual is a race over $V\setminus\{x\}$ that needs no partition bound. SGLang's target-only rule (accept sibling $j$ if $UZ<\sum_{l\le j}w_{x_l}$; otherwise sample the target restricted to untried tokens) is certified by the guards with a cumulative numerator and an exclusion race.
\end{theorem}
\begin{proof}
Positive-part monotonicity gives the enclosures; the rest is the proof of Theorem~\ref{thm:residualrace}. For the point mass, $w_x-Zq_x=w_x-Z\le0$.
\end{proof}
The engine samples that residual by inverse CDF, so a race matches it only in distribution. For a support chosen after the draft pass, retaining the top $K+1$ perturbed entries per tile tightens the outside-support maximum after excluding up to $K$ support tokens; it is not needed for validity, since $r_i\le w_i$ makes the plain tile maximum an upper bound. For a bonus sample the construction is simpler still: a certified race over the target scores with the bonus position's field, with no partition function.

\subsection{A demand frontier over a fixed protocol}\label{sec:frontier}
For a chain of $L$ proposals each head row is accepted, rejected or unresolved. The first certified rejection at position $r$ proves that no row after $r$ can affect the cycle's output, while earlier unresolved rows still matter because one of them may be the true first rejection.

\begin{proposition}[Safe suppression after a certified rejection]\label{prop:frontier}
If every issued certificate equals the decision of a fixed reference verifier on its row, the first true rejection is no later than the first certified rejection. Suppressing head work strictly after the latter does not change the reference output, provided completion and state commit use the actual first rejection.
\end{proposition}
\begin{proof}
At the certified rejection the reference decision is rejection; an earlier unresolved row may also reject, and no later row can be the first rejection. All accepted outputs, the replacement token and the state prefix are determined by rows up to the first rejection.
\end{proof}
The statement is for chains. SGLang's trees share one coin across siblings and choose the next coin by the accepted node, so a tree frontier must follow the tree's paths. The target backbone has already run over every admitted position, so the frontier saves head work, not transformer layers.

\subsection{A cheap decision is not a cheap verification}\label{sec:synthetic}\label{app:synthetic}
The inherited CPU experiment makes the last point concrete. It constructs two families of $128\times8$ integer heads with 16-row tiles and 36 paired fixtures each: one clusters rows around separated centres, the control uses unstructured rows. For each fixture $h^t=h^d+d\,v$ with drift scale $d\in\{0,1,2,4\}$; proposals are sampled from the anchor head's law, and all probability arithmetic is exact rational. The experiment counts distinct target rows projected, which is work, not time.

At zero drift the summaries describe the target exactly and the acceptance decision needs only its candidate row, an identity check rather than a forecast. At nonzero drift the unstructured head defeats greedy screening. More subtly, large drift makes rejection easy to certify, which lowers the rows needed for the decision while destroying acceptance itself, and completing the residual then erases the saving (Figure~\ref{fig:synthetic}). At drift four the clustered family decides with 11.4\% of the rows but needs 96.5\% once the dense residual is charged. The experiment does not refute transport on trained heads; it refutes the inference that a cheap acceptance decision means cheap verification. The residual races of Appendix~\ref{app:guards} are the answer it motivates, and they remain to be measured; fixed-noise verification (Section~\ref{sec:fixednoise}) avoids the residual altogether.

\begin{figure}[tbp]
\centering
\begin{tikzpicture}
\begin{axis}[width=.6\linewidth,height=5.2cm,xmin=-.2,xmax=4.2,ymin=0,ymax=1.05,xtick={0,1,2,4},
  xlabel={Drift scale $d$},ylabel={Fraction of target rows projected},
  grid=major,grid style={draw=ink!10},legend style={at={(1.02,.5)},anchor=west,draw=none,font=\footnotesize,cells={align=left}},legend cell align=left]
\addplot[thick,teal,mark=*,discard if not={family}{clustered},unbounded coords=discard] table[x=drift,y=accept_projected_fraction,col sep=comma]{../data/synthetic_drift.csv};
\addlegendentry{Clustered: decision}
\addplot[thick,teal,dashed,mark=o,discard if not={family}{clustered},unbounded coords=discard] table[x=drift,y=decision_plus_residual_projected_fraction,col sep=comma]{../data/synthetic_drift.csv};
\addlegendentry{Clustered: decision\\+ dense residual}
\addplot[thick,ochre,mark=square*,discard if not={family}{unstructured},unbounded coords=discard] table[x=drift,y=accept_projected_fraction,col sep=comma]{../data/synthetic_drift.csv};
\addlegendentry{Unstructured: decision}
\addplot[thick,ochre,dashed,mark=square,discard if not={family}{unstructured},unbounded coords=discard] table[x=drift,y=decision_plus_residual_projected_fraction,col sep=comma]{../data/synthetic_drift.csv};
\addlegendentry{Unstructured: decision\\+ dense residual}
\addplot[thin,muted,mark=triangle,discard if not={family}{clustered},unbounded coords=discard] table[x=drift,y=acceptance_fraction,col sep=comma]{../data/synthetic_drift.csv};
\addlegendentry{Clustered: accepted}
\end{axis}
\end{tikzpicture}
\caption{Work counts on constructed heads, not timings \ev{synthetic}. Solid lines count the rows needed to decide acceptance; dashed lines add a dense residual completion after each rejection. The thin line is the fraction of proposals accepted in the clustered family. Evidence production, bounds, bonus sampling and the backbone are excluded.}
\label{fig:synthetic}
\end{figure}

